\chapter{Research Methodology \& Evaluation Protocol (Phương pháp và quy trình đánh giá)}

Chương này cần được thống nhất trước khi benchmark.
Nếu mỗi người tự chọn dataset, metric và cách đo, Part III sẽ không so sánh được.

\section{Dataset}

Dataset được chọn phải có camera đúng track, LiDAR, calibration, pose và thời gian ghi đủ tin cậy.

\begin{table}[H]
  \centering
  \caption{Dataset theo track}
  \begin{tabularx}{\linewidth}{l X X}
    \toprule
    \textbf{Track} & \textbf{Dataset} & \textbf{Sequence} \\
    \midrule
    A: Pinhole + LiDAR & TBD & TBD \\
    B: Fisheye + LiDAR & TBD & TBD \\
    \bottomrule
  \end{tabularx}
\end{table}

Mọi approach trong cùng track chạy trên cùng dataset và cùng sequence.

\section{Camera calibration}

Mọi module đọc thông số camera từ một format chung, gồm loại camera, intrinsics, distortion và vị trí so với xe.
Mỗi sequence phải qua một phép kiểm tra trước khi dùng: chiếu LiDAR lên ảnh và xác nhận hai nguồn khớp nhau.
Mô hình fisheye cụ thể được chốt sau khi chọn dataset: \todo{chốt mô hình fisheye}.

\section{Pose conventions}

Các quy ước dưới đây áp dụng cho dữ liệu trao đổi giữa các module. Method có quy ước riêng thì adapter phải tự chuyển đổi.

\begin{table}[H]
  \centering
  \caption{Quy ước dữ liệu}
  \begin{tabularx}{\linewidth}{l X}
    \toprule
    \textbf{Nội dung} & \textbf{Quy ước} \\
    \midrule
    Đơn vị khoảng cách & Mét (đề xuất) \\
    Mốc thời gian & TBD \\
    Hệ tọa độ thế giới & TBD \\
    Hệ tọa độ camera & TBD \\
    Biểu diễn pose & TBD \\
    Nguồn pose & TBD, ưu tiên pose có sẵn nếu kiểm tra thấy đáng tin \\
    \bottomrule
  \end{tabularx}
\end{table}

\section{Experiment conditions}

Để so sánh công bằng, các điều kiện sau phải giống nhau giữa các approach trong cùng track:
\begin{itemize}
  \item Cùng sequence và cùng danh sách ảnh train, ảnh giữ lại để test.
  \item Cùng độ phân giải đánh giá.
  \item Cùng baseline đối chứng: TBD.
  \item Mọi thay đổi cấu hình so với mặc định của repository đều được ghi lại.
\end{itemize}

\subsection{Vertical slice}

Cần có sớm một luồng chạy tối thiểu xuyên qua nhiều thành phần.
Ví dụ: dữ liệu KITTI-360 $\to$ calibration/poses $\to$ một reconstruction method $\to$ render held-out view $\to$ tính PSNR và ghi runtime.

Song song, có thể chạy CARLA với bản đồ có sẵn để kiểm tra simulation interface độc lập. Hai phần này chưa cần nối ngay.

\section{Metrics}

\begin{table}[H]
  \centering
  \caption{Các chỉ số đánh giá}
  \begin{tabularx}{\linewidth}{>{\raggedright\arraybackslash}p{3.2cm} l X}
    \toprule
    \textbf{Nhóm} & \textbf{Chỉ số} & \textbf{Ý nghĩa} \\
    \midrule
    Chất lượng ảnh & PSNR & Sai khác theo từng pixel, càng cao càng tốt \\
     & SSIM & Giống nhau về cấu trúc ảnh, càng cao càng tốt \\
     & LPIPS & Khác biệt theo cảm nhận thị giác, càng thấp càng tốt \\
    \addlinespace
    Hình học & TBD & Sai lệch so với LiDAR \\
    \addlinespace
    Chi phí & Runtime, VRAM & Thời gian chạy và bộ nhớ GPU lớn nhất \\
     & Thao tác thủ công & Số việc phải làm tay cho mỗi scene \\
    \addlinespace
    Determinism & $\Delta p(t)$ & Lệch vị trí xe giữa 2 lần chạy tại cùng thời điểm \\
    \bottomrule
  \end{tabularx}
\end{table}

Với hai lần chạy cùng kịch bản, độ lệch vị trí được tính như sau:
\begin{equation}
  \Delta p(t) = \lVert p_{\text{lần 1}}(t) - p_{\text{lần 2}}(t) \rVert
\end{equation}

\section{Hardware measurement}

Mỗi kết quả ghi kèm GPU, RAM, phiên bản driver và thời điểm chạy.
Các số đo so sánh với nhau nên được chạy trên cùng loại máy.
Nếu buộc phải dùng máy khác, cần ghi rõ để không so sánh trực tiếp thời gian chạy.
Cấu hình máy của nhóm: \todo{bổ sung GPU, VRAM, dung lượng lưu trữ}.

\section{Quy tắc ghi kết quả}

\begin{itemize}
  \item Mỗi kết quả gắn với experiment ID, config và code version để chạy lại được. Experiment ID có dạng \texttt{approach-track-sequence-số}, ví dụ \texttt{a-pinhole-seq01-001}.
  \item Thực nghiệm chưa chạy ghi \texttt{Not Evaluated}, không điền số liệu mẫu vào bảng chính.
  \item Chức năng approach không có ghi \texttt{Not Supported}.
\end{itemize}

\section{Phiên bản protocol}

Protocol có thể cập nhật khi Phase 2 phát hiện vấn đề, nhưng mọi thay đổi phải được ghi lại.

\begin{table}[H]
  \centering
  \caption{Lịch sử thay đổi protocol}
  \begin{tabularx}{\linewidth}{l l X}
    \toprule
    \textbf{Phiên bản} & \textbf{Ngày} & \textbf{Thay đổi} \\
    \midrule
     &  &  \\
    \bottomrule
  \end{tabularx}
\end{table}
